\section{An explicit threshold for independent weights}\label{sec:general}
Fix \(k\ge3\), put \(d=k-1\), and let \(w_1=0\), \(w_a\in[0,1]\) for
\(2\le a\le d\), and \(w_a=1\) for \(a>d\).
In this section \(F_j\) denotes the weighted strict-excedance row from
Theorem~\ref{general}. Use \(p=n-1\), \(r=e^{-x}\), and \(t=-r\), and retain
\(H_j,Y_j,V_j,\rho,c,s,Z\) from Sections~\ref{transport} and~\ref{sec:normalized}.
Throughout the section assume \(n\ge10^6d^2\). We prove \eqref{goal}
on the intervals cut at \(12d+100\) and \(3\sqrt p\).

\subsection{The normalized generating function and its insertion identity}
Define \(A_0(t)=1\) and
\[
 M_a=\frac{t A_{a-1}(t)}{(1-t)^a},\qquad
 b_a=\frac{(w_a-1)M_a}{a!},\qquad
 B(z)=\exp\!\left(\sum_{a=1}^{d} b_a z^a\right).
\]
The labelled cycle formula gives \(\sum H_j z^j/j!=(1-t)B(z)L(z)\).
For \(a\ge2\) the cycle polynomial is \(t A_{a-1}(t)\). For the linear
term, rewriting the unrestricted generating function in terms of \(L\)
produces an extra \(e^z\); deleting the fixed points multiplies this by
\(e^{-z/(1-t)}\), giving \(e^{b_1z}\).
The Eulerian recurrence gives \(M'_a=-M_{a+1}\), with differentiation
in \(x\). In particular \(b_1=-(\log(1-t))'\), and
\(b'_a+(a+1)b_{a+1}=(w_{a+1}-w_a)M_{a+1}/a!\), taking \(b_{d+1}=0\).
Since \((\partial_x+\partial_z)L=0\), coefficient extraction gives
\begin{align}
 C_a&=(w_{a+1}-w_a)M_{a+1}/a!,\nonumber\\
 H_n+H'_p&=\sum_{a=1}^{d}(p)_a C_a H_{p-a},
 &Z&=\frac1n\sum_{a=1}^{d} C_a\rho^{a+1}Y_{p-a},\label{general-insertion}
\end{align}
where \((p)_a=p!/(p-a)!\). The constant term of
\((\partial_x+\partial_z)\log((1-t)B)\) vanishes.
We use \(|Z|<|Y_n|\) to prove \eqref{goal} on the first two intervals.

Write \(v=V(\rho)\), \(V(R)=\sum_{a=1}^d |b_a|R^a\) and
\(J=\sum_{a=1}^d |C_a|\rho^{a+1}\). For every \(R\ge0\) the absolute
coefficient mass of the tilt satisfies
\[
 \sum_{i\ge0}|[z^i]B|R^i\le e^{V(R)} .
\]
Since the exponent of \(B\) has degree at most \(d\), each term in
\(V((1+1/d)\rho)\) is at most \((1+1/d)^d<3\) times the
corresponding term in \(V(\rho)\). Thus
\begin{equation}\label{quadratic-dilation}
 V((1+1/d)\rho)\le 3v .
\end{equation}

\subsection{Two bounds for the tilt}
For all \(x\ge0\) the pole formula in Section~4 gives, when \(a\ge2\),
\[
 \frac{\rho^a|M_a|}{(a-1)!}
 =|V_{a-1}|\le 2+\frac1{as^2}.
\]
The equality follows by the expansion of \(L\) at zero for \(x>0\)
and by continuity at \(x=0\). For the bound, write \(\ell=2b+1\ge3\);
then \(|c+\ii\ell s|^{-a}\le(2ab(b+1)s^2)^{-1}\) by Bernoulli's
inequality, and \(\sum_{b\ge1}1/(b(b+1))=1\). In particular this
bound is below 7 if \(x^2\le32a\), since \(\pi^2>9\). The bound
applies also at order \(a+1\) under the same condition. We will use
\begin{equation}\label{quadratic-tilt-cost}
 v\le11(1+\log d),\qquad J\le24d,\qquad
 x^2>32d\ \Longrightarrow\ v\le1+12d/x.
\end{equation}
For \(x\le8\) use the bound 7 for every nonlinear tilt order and
every insertion order. The linear tilt costs at most
\(\rho e^{-x}\le(x+\pi)e^{-x}<4\). This gives the first two estimates.
The last hypothesis is impossible when \(x\le8\).

For \(x\ge8\) put \(A=\min(d,\lfloor x^2/32\rfloor)\). Orders
\(a>A\) cost at most \(7\sum_{a>A}1/a\) in \(v\) and \(7(d-A)\)
in \(J\). For the low orders use
\(M_a=\sum_{h\ge1}(-r)^h h^{a-1}\), absolutely convergent for \(x>0\).
Their costs are bounded by
\[
 \sum_{h\ge1}\frac{e^{-hx}}h\sum_{a=1}^A\frac{(h\rho)^a}{a!},
 \qquad
 \rho\sum_{h\ge1} e^{-hx}\sum_{a=1}^A\frac{(h\rho)^a}{a!},
\]
respectively. Here \(x\le\rho\le9x/8\) and
\(\rho-x=\pi^2/(\rho+x)\le 8/x\). For \(h\le4A/x\), complete the
exponential and use \(h(\rho-x)\le32A/x^2\le1\). The head of
the first sum is at most \(3(1+\log d)\) and also at most
\(12A/x\). The head of the second is at most \(12\rho A/x\).
For \(h>4A/x\), the truncated exponential is at most
\(2^A e^{h\rho/2}\). Since \(\log 2<3/4\), inclusion of \(e^{-hx}\)
leaves a bound of \(e^{-hx/5}\).
The tails cost at most \(5/x\) and \(5\rho/x\), respectively.
These estimates give the first two bounds of \eqref{quadratic-tilt-cost},
using \(12\rho A/x+5\rho/x\le14A+6\). When \(x^2>32d\) one has
\(A=d\), and the bound \(1+12d/x\) follows from the low orders alone.

\subsection{A centered estimate in the compact interval}
Write \(z_\ell=(c+\ii\ell s)^{-1}\) for odd signed \(\ell\), and
\(D=|B(\zeta)|\), where \(\zeta=x+\ii\pi\). In particular
\(D\ge e^{-v}>0\). For \(\nu\ge8\),
\begin{equation}\label{quadratic-pole}
 \sum_{b\ge1}(2b+1)|z_{2b+1}|^\nu
 \le\frac{4}{(\nu s^2)^2}e^{-\nu s^2/2}.
\end{equation}
In fact, extract half the power from each summand. This part is at most
\((1+8s^2)^{-\nu/4}\le e^{-\nu s^2/2}\), by concavity of
\(u\mapsto\log(1+8u)\) on \([0,1]\) and \(\log9>2\).
For the other half apply
\((1+4b(b+1)s^2)^{\nu/8}\ge(\nu/2)b(b+1)s^2\) twice.
The resulting bound telescopes since
\((2b+1)/(b^2(b+1)^2)=b^{-2}-(b+1)^{-2}\).
Also \(|z_\ell-c|\le|\ell|s\), by the identity
\(z_\ell-c=s(s-\ii c\ell)/(c+\ii\ell s)\).

Put \(\Pi_j=2\Re(B(\zeta)z_1^{j+1})\).
We claim that
\begin{equation}\label{quadratic-centered}
 |Y_j-\Pi_j|< D/8\quad(n-d-1\le j\le n),\qquad
 |(Y_n-cY_p)-(\Pi_n-c\Pi_p)|<sD/8.
\end{equation}
To see this, write \(B_i=[z^i]B\).
The exact row is the convolution \(Y_j=\sum_{i=0}^{j}B_i\rho^i V_{j-i}\).
The principal term is the corresponding sum over all \(i\ge0\)
using the nearest pair, since \(z_1^{-1}=\zeta/\rho\).
Use the same head \(i\le T=\lfloor n/2\rfloor\) in both rows of the
centered difference. Each pole exponent in the ordinary head is
\(\nu=j-i+1\ge n/3\), since \(n\ge6d\); in the centered head it is \(\nu=n-i\ge n/3\),
with additional factor exactly \(z_\ell-c\).
In this interval \(\rho\le64d\), so \(\nu s^2\ge2\).
The ordinary head divided by \(D\), and the centered head divided by
\(sD\), are each at most \(2\exp(2v-ns^2/6)\), by
\eqref{quadratic-pole} and its conjugate. This is less than \(1/16\).
If \(x^2\le32d\) then \(\rho^2\le40d\), so
\(ns^2/6\ge (9\cdot10^6/240)d\), while \(2v\le22d\).
Otherwise write \(y=d/x\ge1/62\).
Then \(\rho^2\le2x^2\), giving \(ns^2/6\ge750000y^2\), and
\(2v\le2+24y\). The difference exceeds \(\log32\) at \(y=1/62\)
and is increasing thereafter.

For the tail, \(|V_j|\le\rho^2\) at all nonnegative orders, by the
pole bound above or the formula for \(V_0\). Every two-pole term
has modulus at most two before multiplication by \(B_i\rho^i\).
With \(R=(1+1/d)\rho\) we have
\[
 \sum_{i>T}|B_i|\rho^i\le e^{V(R)}(1+1/d)^{-(T+1)}.
\]
This includes the principal terms beyond the finite row convolution.
Therefore either tail divided by its respective scale has bound
\[
 \frac{4\rho^2}{s}\exp\!\left(V(R)+v-\frac{n}{2(d+1)}\right)<1/16.
\]
Indeed \(4\rho^2/s<350000 d^3\) and by \eqref{quadratic-dilation}
and \eqref{quadratic-tilt-cost} the positive exponent including the
prefactor and \(\log16\) is less than
\(23+3\log d+44(1+\log d)\le57d\).
But \(n/(2(d+1))\ge10^6d/3\).
This proves \eqref{quadratic-centered}.

At \(Y_p=0\), put \(S=B(\zeta) z_1^{p+1}\).
Then \(|\Re S|<D/16\), so \(|\Im S|>15D/16\). Since
\(\Pi_n-c\Pi_p=2s\Im S\), the centered estimate yields
\(|Y_n|>7sD/4\). The shorter rows in \eqref{general-insertion}
have magnitudes below \(17D/8\). From \(\rho\le64d\) and
\(J\le24d\) one gets \(J/s\le512d^2\). Thus \(|Z|<|Y_n|\).

\subsection{The intermediate interval}
On \(12d+100\le x\le3\sqrt p\), we have \(\rho\le4x/3\) and
\(ns^2\ge1\). In the positive-label sums above take \(A=d\) and
bound the whole truncated exponential by \(2^d e^{h\rho/2}\).
This yields
\[
 v\le2^{d+1}e^{-x/3},\qquad J\le\rho\,2^{d+1}e^{-x/3},\qquad
 \rho^3 2^{d+1}e^{-x/3}<1/32 .
\]
For the last bound use \(d\le(x-100)/12\), \(\log2<1\), and the
decrease of \(x^3e^{-x/4}\) on \([100,\infty)\).
The upper bound \(2(4/3)^3 100^3/(2^8 2^{25})\) is below \(1/32\).
Now \(|V_j|\le\rho^2/2\) by the telescoping bound of Section~4,
including order zero. The absolute tilt majorant gives
\[
 |Y_j-V_j|\le (\rho^2/2)(e^v-1)\le\rho^2v<s/32 .
\]
At \(Y_p=0\) the bare phase gap \eqref{center} gives
\(|V_n-cV_p|>s\), hence \(|Y_n|>15s/16\).
Since \(|Y_j|\le\rho^2\), the insertion is at most
\(\rho^2J/n<1/(32n)\le s/32\). This proves \eqref{goal}.

\subsection{The exterior from the exact positive labels}
For \(x\ge3\sqrt p\) we use \(p\ge\max(65536,4k^2)\), which
follows from \(n\ge10^6d^2\) and \(k=d+1\le2d\).
Write \(u=\sqrt p\). For a positive integer \(h\), choose a uniform
permutation of \(h\) slots and assign \(j\) labelled balls to these
slots independently and uniformly. If \(A_C\) is the number of
balls in a slot cycle \(C\), put \(\eta_j(h)=\mathbb{E}\prod_C w_{A_C}\),
with \(w_0=w_1=0\). The actual coefficient of \(t^h\) in
\(F_j(t)/(1-t)^j\) is \(h^j\eta_j(h)\). Indeed the cycle formula and
the positive-label expression for \(M_a\) write this row EGF as
\[
 \exp\!\left(\sum_{h\ge1}\frac{t^h}{h}
                   \sum_{a\ge2}w_a \frac{(hz)^a}{a!}\right).
\]
The coefficient of \(t^h\) is the cycle-index average over permutations
of the \(h\) slots. For each cycle of length \(b\), the factor
\((bz)^a/a!\) assigns each of the chosen \(a\) balls to one of its \(b\) slots;
\(j![z^j]\) completes the assignments of labelled balls.
At \(h=0\) the coefficient vanishes for \(j>0\).
In particular, with \(S_q=\sum_{h\ge1}(-r)^h h^q\),
\[
 H_j=\sum_{h\ge1}(-r)^h h^j\eta_j(h),\qquad
 -H'_p=\sum_{h\ge1}(-r)^h h^{p+1}\eta_p(h).
\]
Absolute local uniform convergence for \(x>0\) justifies differentiation.

Let \(L=\lfloor u/2\rfloor\). For \(h\le L\) and \(j\ge p\),
the mean \(\mu=j/h\) for each slot is at least \(2u\ge4k\).
When every slot has at least \(k\) balls, all its cycle weights are
one. For any individual binomial slot count \(X\),
Markov's inequality applied to \(2^{-X}\) bounds
\(\Pr(X\le\mu/4)\) by
\(\exp(-\mu/2+\mu\log(2)/4)\le e^{-\mu/4}\).
Thus
\begin{equation}\label{quadratic-slot-head}
 0\le1-\eta_j(h)\le h e^{-j/(4h)}
 \le(u/2)e^{-u/2}<1/(100n).
\end{equation}
We will also use \(e^{u/8}>200u^3\) for \(u\ge256\). It holds
at 256 since \(2^{32}>200\cdot256^3\), and
\(e^{u/8}/u^3\) is increasing there.

For \(q=p,n\), let \(m_q\) be a mode of \(h^q r^h\) over positive
\(h\) and let \(W_q=m_q^q r^{m_q}\).
Comparison with the preceding label gives \(m_q\le q/x+1\le u/3+2\).
For either exterior tail beyond the adjacent labels the weight at
distance \(b\ge2\) from the mode is at most
\(W_q 50^{1-b}\). To see this, the first such outward ratio is at
most \((1-(m_q+1)^{-2})^q\), by taking its ratio with the previous
one, which is at most one. On the left the same bound follows
with \(m_q-1\) in place of \(m_q+1\), whenever that tail exists.
Since \(m_q+1\le3u/7\), both bounds are at most
\(e^{-49/9}<1/50\); subsequent ratios decrease. Consequently
\[
 \sum_{h\ge1}h^q r^h<4W_q,\qquad
 \sum_{h>L} h^q r^h
 <2W_q 50^{3-u/6}\le2W_q e^{-u/8}<W_q/(100n).
\]
Here \(L-m_q\ge u/6-3\ge u/8\).
Combining \eqref{quadratic-slot-head} and the tail, using
\(0\le \eta_j(h)\le1\) everywhere and \(W_n\le nW_p\), gives
\begin{equation}\label{quadratic-exterior-errors}
 |H_p-S_p|\le W_p/(20n),\qquad
 |H_n-S_n|\le W_p/20,\qquad
 |-H'_p-S_n|\le W_p/20.
\end{equation}

Finally take \(h_0\) to be the stronger neighbor of \(m_p\),
with missing left neighbor given weight zero. The smaller neighbor
has weight below \(W_p/10\), since the product of the neighbor
weights at \(m_p>1\) is
\(W_p^2(1-m_p^{-2})^p\) and \(m_p\le 2u/5\).
Centering \(S_n-h_0S_p\) cancels the stronger neighbor and leaves
the mode with magnitude \(W_p\). The remaining immediate term is
bounded by \(2W_p/10\); the more distant terms cost at most
\(2W_p\sum_{b\ge2}(b+1)50^{1-b}\). Their total is strictly below
\(W_p/2\). It follows that
\[
 W_p/2<|S_n|+h_0|S_p|.
\]
At \(H_p=0\), \(h_0\le n\) and \eqref{quadratic-exterior-errors}
give \(|S_n|>9W_p/20\). Both true quantities in \eqref{goal}
have the sign of \(S_n\), proving the assertion in the exterior.

\subsection{The original rows}
For \(j\ge k\), the linear coefficient of the original row is one:
a single \(j\)-cycle has unit weight and its cycle polynomial is
\(tA_{j-1}(t)\), while two or more cycles give at least two excedances.
Inversion gives reciprocity of the quotient, which is monic with
constant term one and has no nonnegative real root.
For the chosen profile connect its short-cycle weights by a straight
segment to the all-one short-cycle vector, keeping the fixed-point
and longer-cycle weights fixed. The sign \eqref{goal} holds over the
whole segment at every predecessor zero in logarithmic coordinates.
Reciprocity and the transport argument of Section~\ref{transport}
give strict interlacing with simple negative quotient roots.
Restoring the common simple zero gives weak interlacing, and the
all-zero short-cycle vector yields the minimum-cycle-length
enumerators. This proves Theorem~\ref{general}.
